\documentclass[11pt]{article}
\usepackage[margin=1in]{geometry}
\usepackage{amsmath,amssymb,amsthm}
\usepackage{mathtools}
\usepackage{booktabs}
\usepackage[colorlinks=true,linkcolor=blue,citecolor=blue]{hyperref}
\usepackage[capitalise,nameinlink]{cleveref}

\newtheorem{theorem}{Theorem}
\newtheorem{proposition}[theorem]{Proposition}
\newtheorem{lemma}[theorem]{Lemma}
\newtheorem{corollary}[theorem]{Corollary}
\newtheorem{definition}[theorem]{Definition}
\newtheorem{remark}[theorem]{Remark}

\newcommand{\Lean}{Lean}
\newcommand{\code}[1]{\texttt{#1}}
\newcommand{\claim}[1]{\textbf{#1}}
\newcommand{\RR}{\mathbb{R}}
\newcommand{\NetG}{\mathrm{Net}}
\newcommand{\VecG}{\mathrm{Vec}}
\newcommand{\LinOf}{\mathrm{linOf}}
\newcommand{\PatOf}{\mathrm{patternOf}}
\newcommand{\RegOf}{\mathrm{regionOf}}

\emergencystretch=1.5em

\title{Exact explanations for piecewise-linear networks are\\
unique and local: a machine-checked account}
\author{Fan Li\\[8pt]
{\small ORCID: 0009-0004-8329-3642}\\[2pt]
{\small \texttt{56213760@qq.com}}}
\date{}

\begin{document}
\maketitle

\begin{abstract}
Post-hoc explanation methods for neural networks --- saliency maps, Grad-CAM,
Shapley values --- are heuristics, and the question usually asked of them is
whether they are faithful. We argue that the question should be replaced by a
geometric one: for a piecewise-linear network there is an exact local model,
and one should ask what that model is, whether it is unique, how far from a
point it remains valid, and what determines its shape.

This paper gives machine-checked answers for a class of convolutional
networks. We formalize networks as finite compositions of affine maps,
pointwise ReLU and max-pooling, indexed by an arbitrary finite additive
commutative semigroup, so that one- and two-dimensional circular convolution
are instances of the same theorem rather than two theories. We prove that every
input has an activation pattern whose cell contains it and on which the network
is exactly the pattern's linearization; that within an explicit neighborhood ---
given by per-layer margins and pooling gaps --- the linearization remains valid;
that the linearization is the \emph{unique} affine map agreeing with the
network on that neighborhood; and that its zero pattern is determined by the
architecture. We also prove a negative result: with zero padding instead of
circular boundary, translation equivariance fails, and we exhibit the failure.

Everything is formalized in \Lean{}~4 with Mathlib across twenty ledger entries,
each with an independently written challenge statement checked against the
declaration's type. We are explicit about what the formalization does not
establish: no compactness, no explicit region geometry, and a model class that
is stylized rather than a production CNN.

\end{abstract}

\tableofcontents

\section{Introduction}

\subsection{The question that explanation methods do not answer}

A great deal of work on interpretability asks whether a particular explanation
method is \emph{faithful}: whether a saliency map reflects the computations
that produced a prediction. The answers are necessarily method-by-method, and
typically empirical.

For piecewise-linear networks there is a sharper object available. Such a
network is affine on each cell of a finite partition of its input space, so at
every input there is an \emph{exact} local model: an affine map that agrees
with the network on a neighborhood of that input. The interesting questions are
then not about any particular method but about that model:

\begin{enumerate}
\item Does it exist? (Yes, and its cell contains the point.)
\item Is it unique? (Yes, on the cell's interior --- which needs proof.)
\item How far is it valid? (An explicit radius, from margins and pooling gaps.)
\item What determines its shape? (The architecture, through the receptive
  field.)
\end{enumerate}

These are mathematical questions with mathematical answers, and this paper
gives those answers for a family of convolutional networks, entirely inside a
proof assistant.

\subsection{What is proved}

The ledger has twenty entries; we highlight six.

\begin{center}
\begin{tabular}{@{}p{0.07\textwidth}p{0.83\textwidth}@{}}
\toprule
Ledger & Statement \\
\midrule
C5 & For every network and input $x$ there is a pattern $\sigma$ with
  $x \in \RegOf(\sigma)$ and $f(y) = \LinOf(\sigma)(y)$ for all $y \in \RegOf(\sigma)$. \\
C6 & In particular $f(x) = \LinOf(x)(x)$: the prediction is determined by the
  activation pattern. \\
C17 & If the margins and pooling gaps at $x$ satisfy a recursive condition
  $\mathrm{Bounds}$, then the ball of radius $R$ around $x$ lies inside
  $\RegOf(\sigma)$ --- the linearization is valid on the whole ball. \\
C18 & The linear map's nonzero entries lie in a recursively defined receptive
  field: the zero pattern is architectural, not data-dependent. \\
C20 & If $\mathrm{Bounds}$ holds with $R>0$, every affine map agreeing with
  the network on that ball \emph{equals} $\LinOf(x)$. The explanation is unique. \\
C13 & With zero padding instead of a circular boundary, translation
  equivariance fails; an explicit three-element counterexample is given. \\
\bottomrule
\end{tabular}
\end{center}

C5 and C6 are the familiar statement that a ReLU network is piecewise linear.
We do not claim them as new mathematics, and
\cref{sec:limits} says so plainly. C17, C18 and C20 are the parts that a
statement of the familiar fact does not give: a validity radius, an
architectural determination of the linear map's sparsity, and uniqueness.

\subsection{What is not claimed}

\begin{itemize}
\item \textbf{No compactness.} The number of cells grows exponentially with
  depth; an exact explanation of size polynomial in the network does not exist
  in general. We prove nothing in that direction and say nothing about
  ``human-readable'' explanations.
\item \textbf{No production CNNs.} The formalized class is stylized: equal
  width at every layer, no downsampling, no flattening or classifier head,
  bias-free convolution, and only ReLU and max-pooling as nonlinearities.
  Section~\ref{sec:limits} states this in full.
\item \textbf{No claim of novelty for the piecewise-linear fact itself}, which
  is classical.
\end{itemize}

\subsection{Formalization}

All results are in \Lean{}~4 with Mathlib. Each ledger row with status
\claim{PROVED} has a declaration whose type has been compared, by the
elaborator, against a statement written separately in a challenge file; and
each has been axiom-audited, with the transitive axiom set contained in
$\{\code{propext}, \code{Classical.choice}, \code{Quot.sound}\}$. A kernel
replay of the compiled artifacts through \code{leanchecker} also accepts the
development.

The development was audited by three independent reviewers with different
lenses (statement fidelity, proof correctness, mathematical content). One
substantive defect was found and is recorded in \cref{sec:earlier}: a claim in
an early draft of the ledger that the explanation ``decomposes by receptive
field'' had no theorem behind it. The theorem now exists (C18).

\section{The model}

\subsection{Index sets}

Fix a finite additive commutative semigroup $G$ with decidable equality. The
primary instances are $G = \mathbb{Z}/n$ (written \code{Fin n}) for
one-dimensional circular convolution and $G = \mathbb{Z}/h \times \mathbb{Z}/w$
for two dimensions. We use \code{AddCommSemigroup} rather than
\code{AddCommMonoid} deliberately: \code{Fin n} is an additive commutative
semigroup for every $n$, but a monoid only when $n \ne 0$, and the tensor
statements should not require a nonzero size.

A \emph{vector} is a function $G \to \RR$, written $\VecG G$; a \emph{matrix} is
a $\VecG G$-indexed square matrix, written $\mathrm{Mat}\,G$.

\subsection{Layers and networks}

A layer is one of three things:
\begin{itemize}
\item an affine map $x \mapsto Ax + b$, with $A$ a matrix and $b$ a vector;
\item a pointwise ReLU, $x \mapsto (\max(x_i, 0))_i$;
\item a max-pooling layer, given a family of nonempty windows
  $(w_i)_{i \in G}$ with $w_i \subseteq G$, sending $x$ to
  $({\max_{j \in w_i} x_j})_i$.
\end{itemize}
A network is a finite list of layers, evaluated by composition.

Convolution is not a separate constructor: it is a particular affine layer,
whose matrix is circulant (Section~\ref{sec:structure}). This is the design
decision that keeps the development to one theory.

\section{Exact linearization}

\subsection{Patterns and their regions}

For a vector $x$, the ReLU \emph{pattern} records which coordinates are
positive:
\[
  \pi(x)_i \;=\; [\,0 < x_i\,] \in \{\mathrm{true}, \mathrm{false}\}.
\]
The branch convention matters: $x_i = 0$ is recorded as $\mathrm{false}$, and
since $\max(0,0) = 0$ agrees with the closed branch, the linearization below is
exact on the \emph{closed} cell, not merely its interior (ledger C7).

A pooling layer contributes a different kind of pattern: the coordinate at
which each window attains its maximum. Since ties may occur, we fix a choice
function; the value of the explanation does not depend on which tied coordinate
is chosen, only the representation does.

A pattern of a network is therefore a list of items, each either a ReLU switch
vector or a pooling choice; its length equals the number of ReLU and pooling
layers.

\subsection{A forward recursion}

The linearization must be computed \emph{forward}. The reason is structural:
the pattern of a ReLU layer depends on that layer's own input, which is the
prefix of the network applied to $x$. One cannot decide the patterns first and
compute the prefix afterwards.

The recursion \code{walk} carries three things: the current value, the affine
map from the original input to the current value, and the patterns collected so
far. An affine layer composes on the way back without adding a pattern; a ReLU
layer adds its switch pattern and composes with the corresponding diagonal
$0/1$ matrix; a pooling layer adds its choice and composes with the
corresponding selection matrix.

From \code{walk} we read off $\LinOf(\mathrm{net}, x)$, the affine map, and
$\PatOf(\mathrm{net}, x)$, the pattern list. The region is the fibre of the
pattern map,
\[
  \RegOf(\sigma) \;=\; \{\, y : \PatOf(\mathrm{net}, y) = \sigma \,\}.
\]

\subsection{The theorem}

\begin{theorem}[Ledger C5]
For every network $\mathrm{net}$ and input $x$,
\[
  x \in \RegOf(\PatOf(x)) \quad\text{and}\quad
  \forall y \in \RegOf(\PatOf(x)),\ \ f(y) = \LinOf(x)(y).
\]
\end{theorem}

The proof is an induction on the layer list whose hypothesis is equality of
\emph{pattern sequences}, not equality of inputs. The base case is the identity;
the affine step transports the hypothesis along the affine map; the ReLU and
pooling steps use, respectively, that a ReLU layer equals its own
linearization (C3) and that a pooling layer equals the selection matrix at its
own choice (the pooling half of C10).

The corollary is the completeness statement: $f(x) = \LinOf(x)(x)$ (ledger C6).
The prediction is determined by the activation pattern, and the determination
is an explicit affine formula.

\section{Structure}\label{sec:structure}

\subsection{Convolution is circulant}

Index arithmetic in \code{Fin n} is modular. Subtraction modulo $n$ is painful
to formalize, so the convolution is written with addition only:
\[
  (k \star x)_i \;=\; \sum_{m} k_m\, x_{i+m}.
\]
This is cross-correlation in the usual machine-learning terminology; because
$k$ ranges over the whole index set it equals the convolution with the flipped
kernel, so no result depends on the naming, but the paper must use the correct
word.

The associated matrix is $A_{ij} = \sum_m [\,j = i+m\,]\, k_m$, a circulant,
and the bridge $A x = k \star x$ is proved (so the convolution is an affine
layer in the sense of Section 2). The mathematical content of ``circulant'' is
that convolution commutes with translation:
\[
  k \star (x \circ \tau_t) \;=\; (k \star x) \circ \tau_t ,
\]
where $\tau_t$ shifts the index by $t$ (ledger C8).

\subsection{Pooling is selection}

A pooling layer is not affine and not pointwise, so it cannot be an affine
layer. What saves it is that it is a \emph{selection}: at every input the
maximum over each window is attained at some coordinate, and the pooling layer
equals the matrix selecting those coordinates (ledger C9). The choice of tied
coordinate is arbitrary but the value is not.

\subsection{The receptive field}

The linearization of a whole network is a single dense matrix, but its zero
pattern is not arbitrary. Define, backwards through the layers, the set of
input coordinates on which output coordinate $i$ can depend:
\begin{itemize}
\item affine: the union, over the current support, of the nonzero rows of $A$;
\item ReLU: keep only the coordinates whose switch is on;
\item pooling: replace $j$ by the chosen coordinate $c_j$.
\end{itemize}
\begin{theorem}[Ledger C18]
If the linear map's $(i,j)$ entry is nonzero, then $j$ lies in this set.
\end{theorem}
The zero pattern is thus determined by the architecture; the input enters only
through the pooling choices. This is what ``the explanation decomposes by
receptive field'' should mean, and it is now a theorem rather than a phrase.

\subsection{Two dimensions}

Because the index type is an arbitrary finite additive commutative semigroup,
two-dimensional circular convolution requires no new proof: it is the same
theorem at $G = \mathbb{Z}/h \times \mathbb{Z}/w$ (ledger C11). The convolution
and the network both instantiate, and the equivariance statement likewise
(C12). One theorem, two instances.

\section{The boundary}

Circular convolution has no boundary: index addition wraps, translation is a
bijection, and equivariance is unconditional. Zero padding introduces a
boundary, and the price is exactly local equivariance. Rather than state a
weakened positive result, we state and prove the failure, with a witness
(ledger C13): with $n = 3$, a kernel keeping only offset $1$, input
$[1,2,3]$ and a shift by $1$,
\[
  \operatorname{conv}(K, \text{shift}\,x) = [3,1,0], \qquad
  \text{shift}(\operatorname{conv}(K,x)) = [3,0,2],
\]
which differ at coordinate $1$. The reason is local: at output coordinate $2$
the offset leaves the grid and is filled with $0$, while the shifted input sees
a value there.

\section{Local validity}\label{sec:radius}

Existence of an explanation at a point does not say how far from that point the
explanation can be trusted. This section gives an explicit neighborhood.

\subsection{Stability of switches and choices}

Two elementary facts.

\begin{lemma}[Ledger C14]
If every coordinate of $x$ is at distance at least $\varepsilon$ from $0$, and
$y$ differs from $x$ coordinatewise by at most $\varepsilon$, then $y$ and $x$
have the same ReLU pattern.
\end{lemma}

\begin{lemma}[Ledger C15]
If, at $x$, every pooling window's maximum exceeds its other entries by at
least $\delta$, and $y$ differs from $x$ coordinatewise by at most
$\delta/2$, then $y$ and $x$ have the same pooling choices.
\end{lemma}

The strictness in C15 is not cosmetic: the choice is an arbitrary function of
the input, so under a tie a different input may choose differently. A strict
gap rules out ties on the whole neighborhood.

Two conventions deserve comment. We write closeness \emph{non-strictly}
($|y_i - x_i| \le r$), because the Lipschitz estimate for a maximum is
non-strict and the inequality propagates through pooling; the margins and gaps
are correspondingly strict. And the affine step of the recursion multiplies the
radius by $\max(\|A\|_\infty, 1)$ rather than $\|A\|_\infty$, purely to keep
the radius positive when $A = 0$; this is conservative, not tight.

\subsection{The recursive radius}

A predicate $\mathrm{Bounds}(\mathrm{net}, x, R)$ collects the per-layer
conditions along the trajectory: at each ReLU layer, $R$ is below every
coordinate's distance from $0$; at each pooling layer, $2R$ is below every
window's gap; at each affine layer, the radius is multiplied by the row-sum
norm.

\begin{theorem}[Ledger C16]
If $R > 0$ and $\mathrm{Bounds}(\mathrm{net}, x, R)$ holds, then any $y$ with
$|y_i - x_i| \le R$ for all $i$ has the same pattern sequence as $x$.
\end{theorem}

\begin{corollary}[Ledger C17]
Under the same hypotheses, the ball of radius $R$ around $x$ is contained in
$\RegOf(\PatOf(x))$; hence the linearization $\LinOf(x)$ agrees with the network on
that whole ball.
\end{corollary}

$\mathrm{Bounds}$ is a recursive predicate rather than a closed-form number on
purpose: a formula would have to hide the per-layer hypotheses that make the
bound true.

\section{Uniqueness}

Existence and validity still leave open whether the linearization is \emph{the}
explanation or merely one of many. It is the former.

\begin{lemma}[Ledger C19]
Two affine maps that agree on a ball of positive radius are equal.
\end{lemma}

The proof is short and classical: compare the constant terms at the centre, and
the $k$-th columns at the centre displaced by $(r/2)e_k$.

\begin{theorem}[Ledger C20]
If $R > 0$ and $\mathrm{Bounds}(\mathrm{net},x,R)$ holds, then every affine map
that agrees with the network on the ball of radius $R$ around $x$ \emph{equals}
$\LinOf(x)$.
\end{theorem}

The proof uses C17 to place the ball inside the region and C19 to compare.
The consequence is a change of status: an explanation is not a representation
one chooses among many, but a property the network determines. Approximate
explanation methods can then be described precisely --- they approximate a
particular object --- rather than compared to one another empirically.

Note that C17, which we first regarded as a refinement, is in fact the
hypothesis of uniqueness: it is what supplies an open set on which two affine
maps can be compared.

\section{What is not established}\label{sec:limits}

\begin{enumerate}
\item \textbf{The piecewise-linear fact is classical.} C5 and C6 restate, in a
  proof assistant, that a ReLU network is affine on each activation cell. The
  proof is an induction over the layer list; there is no new lemma of substance
  in it. Papers on linear regions have long used the fact; what they add is
  counting, which we do not do.
\item \textbf{No compactness.} The number of cells grows exponentially with
  depth, so no exact explanation of polynomial size exists in general. Nothing
  here contradicts that, and nothing here gives a lower bound either.
\item \textbf{No explicit region geometry.} $\RegOf(\sigma)$ is defined as a
  fibre of the pattern map, not as an intersection of half-spaces. It is
  convex and has nonempty interior when $\mathrm{Bounds}$ holds (by C17), but
  we do not state it that way, and we give no description of the region's
  faces.
\item \textbf{A stylized class.} Equal width throughout; no downsampling,
  stride, flattening or classifier head; bias-free convolution; only ReLU and
  max-pooling. For a production CNN these results are vacuous, not false.
\item \textbf{Noncomputable constructions.} The pattern function uses
  \code{decide} on a real inequality, and the pooling choice uses
  \code{Classical.choose}; neither computes. ``Machine-checkable'' holds in the
  kernel sense, not the executable sense.
\end{enumerate}

\section{Related work}

The piecewise-linear structure of ReLU networks and the counting of linear
regions are studied in \cite{montufar,telgarsky,arora}. Explanation methods
form a large literature; \cite{lime,shap,gradcam,ig} are representatives. The
present work does not propose a method; it formalizes the object such methods
approximate.

Formal verification of neural networks has concentrated on robustness
certification and on verifying properties of trained networks, typically with
interval or abstract interpretation rather than in a general proof assistant.
Our contribution is complementary: a specification-level account of what an
explanation \emph{is}, in the same environment used for the rest of
mathematics.

\section{Formalization notes}\label{sec:earlier}

Three points about the process, recorded because they bear on the results.

\paragraph{The statement gate.} Every \claim{PROVED} row has a challenge
declaration in a separate file, and a script checks that the challenge's type
and the theorem's type are definitionally equal. The check caught an error on
its first run: the challenge quantified the index explicitly while the theorem
had it implicit. That is not a mathematical difference, but it is a difference,
and the gate is deliberately stricter than mathematics.

\paragraph{The audit.} Three independent reviewers examined the development
with different lenses. The mathematical-content reviewer found the substantive
defect mentioned in the introduction: the ledger asserted a receptive-field
decomposition that no theorem supported. It also concluded that, at that stage,
the theorem was not sufficient as a paper's main result. The radius (C14--C17),
the receptive field (C18) and uniqueness (C19--C20) were added in response.

\paragraph{Persistent limitation.} The reviewers share the author's model
family. This is a residual limitation of the audit, not a claim of independent
human review.

\appendix

\section{The ledger}

\begin{center}
\begin{tabular}{@{}lll@{}}
\toprule
ID & Content & Shape \\
\midrule
C1 & affine composition & base algebra \\
C2 & affine base case & induction base \\
C3 & ReLU linearization & induction base \\
C4 & induction step & engine \\
C5 & exact linearization & main \\
C6 & completeness corollary & corollary \\
C7 & boundary agreement at $z=0$ & convention \\
C8 & convolution commutes with shifts & structure \\
C9 & pooling is a selection & structure \\
C10 & CNN instance & instance \\
C11 & two-dimensional CNN & instance \\
C12 & two-dimensional equivariance & instance \\
C13 & zero padding breaks equivariance & negative \\
C14 & ReLU switch stability & local \\
C15 & pooling choice stability & local \\
C16 & recursive radius & local \\
C17 & the ball is inside the region & local \\
C18 & receptive field & structure \\
C19 & affine uniqueness on a ball & uniqueness \\
C20 & the explanation is unique & uniqueness \\
\bottomrule
\end{tabular}
\end{center}

\begin{thebibliography}{99}

\bibitem{montufar}
G.~Mont\'ufar, R.~Pascanu, K.~Cho, and Y.~Bengio.
\newblock On the number of linear regions of deep neural networks.
\newblock \emph{Advances in Neural Information Processing Systems}, 2014.
\newblock arXiv:1402.1869.

\bibitem{telgarsky}
M.~Telgarsky.
\newblock Benefits of depth in neural networks.
\newblock \emph{Conference on Learning Theory}, 2016.
\newblock arXiv:1602.04485.

\bibitem{arora}
R.~Arora, A.~Basu, P.~Mianjy, and A.~Mukherjee.
\newblock Understanding deep neural networks with rectified linear units.
\newblock \emph{International Conference on Learning Representations}, 2018.
\newblock arXiv:1611.01491.

\bibitem{lime}
M.~T. Ribeiro, S.~Singh, and C.~Guestrin.
\newblock ``Why should I trust you?'': explaining the predictions of any
  classifier.
\newblock \emph{ACM SIGKDD}, 2016.
\newblock arXiv:1602.04938.

\bibitem{shap}
S.~M. Lundberg and S.-I. Lee.
\newblock A unified approach to interpreting model predictions.
\newblock \emph{Advances in Neural Information Processing Systems}, 2017.
\newblock arXiv:1705.07874.

\bibitem{gradcam}
R.~R. Selvaraju, M.~Cogswell, A.~Das, R.~Vedantam, D.~Parikh, and D.~Batra.
\newblock Grad-CAM: visual explanations from deep networks via gradient-based
  localization.
\newblock \emph{International Conference on Computer Vision}, 2017.
\newblock arXiv:1610.02391. doi:10.1109/ICCV.2017.74.

\bibitem{ig}
M.~Sundararajan, A.~Taly, and Q.~Yan.
\newblock Axiomatic attribution for deep networks.
\newblock \emph{International Conference on Machine Learning}, 2017.
\newblock arXiv:1703.01365.

\end{thebibliography}

\end{document}
